型。首先，通过坐标系定义与状态变量选取，给出了车辆位姿与控制输入的统
一描述；其次，在纯滚动与低速假设下，建立了车辆非完整约束条件，并推导
了双桥转向工况下的状态方程、转弯半径表达式与曲率模型；随后，结合前桥
单独转向、前后桥反向转向和同向转向三种典型模式，分析了不同转向策略对
车辆机动能力的影响；最后，引入速度、转角与曲率约束，并说明了理论模型
在 ROS-Gazebo 环境中的参数映射方式。
本章模型为后续接近控制和仿真实验提供了基础。相较于普通道路车辆模
型，双桥转向模型更能体现机场牵引车在窄空间内调头、贴近和横向微调的特
点；同时，速度、转角和曲率约束也为第三章的路径生成、误差反馈和停车判
定提供了必要限制。










\section{基于飞机辅助感知的自动接近与对接控制方法}
\subsection{感知任务分析与总体方案设计}
在机场飞机自动牵引作业过程中，牵引车能否安全、准确地完成对接，不
仅取决于车辆本体的运动控制能力，还取决于其对目标飞机及作业环境的感知
能力。特别是在接近飞机前起落架并执行最终对接时，系统需要同时解决目标
身份确认、目标位置识别、姿态估计及近距离安全态势判断等问题。因此，感
知系统是实现自动牵引作业闭环控制的关键组成部分[18,29-30]。
结合本文仿真条件，感知方案采用“飞机辅助标识+车端传感器”的结构。
在前起落架附近设置可识别特征后，牵引车可以通过前向视觉和近距测量获得
目标身份、中心偏移和局部方向。这样做不是单独研究标签识别，而是把形状
复杂、易受视角影响的起落架区域转化为可观测的对接参考点，使车辆能够先
确认目标，再逐步修正位置和姿态[29-30]。
与传统依赖人工目视判断的作业方式相比，该方案具有以下特点：其一，
能够为牵引车提供明确的目标身份信息，避免牵引对象识别错误；其二，能够
在接近过程中持续提供目标相对位置与姿态信息，为控制器提供实时反馈；其
三，能够将飞机对接区域的局部环境显式化，使碰撞风险、偏航误差和横向偏
差能够被量化描述。
因此，飞机辅助感知并不是单纯的“加一个二维码”或“加一个标签”，
而是将目标飞机前起落架区域设计为牵引系统可识别、可验证、可跟踪的感知
节点，从而为后续自动控制提供稳定的先验信息和观测依据。
基于上述分析，本文构建的感知总体方案由三部分组成：
第一部分为目标身份确认模块，用于确认当前牵引任务对应的飞机对象是否与
预设目标一致；
第二部分为相对位姿感知模块，用于估计牵引车与飞机前起落架之间的相
对距离、横向偏差和姿态偏差；
第三部分为近距态势判断模块，用于在对接过程中识别安全边界、碰撞风
险与姿态偏离程度，并向控制层输出约束信息。
从系统结构上看，该方案形成了由“飞机端目标标识—牵引车端感知模块
—控制决策节点”构成的信息链路。飞机前起落架不再只是被动的对接对象，
而是成为局部识别信息载体；牵引车通过传感器采集这些信息，再将其转换为
控制层能够使用的状态量。这样一来，第二章中的车辆运动学模型与本章中的







感知输出自然衔接：前者解决牵引车如何运动，后者解决车辆如何获得目标、
误差和停车依据。
\subsection{飞机辅助感知目标与识别信息设计}
为提高前起落架区域的可识别性，本文把辅助感知对象限定在飞机前部的
局部功能区，并在该区域设置结构化标识。该标识不改变飞机运动学特性，也
不参与牵引受力，只承担识别和定位作用；换言之，它相当于给牵引车提供一
个稳定的局部参考基准。
本文在前起落架区域设置的识别码主要包含目标飞机编号、对接区域中心
位置、局部方向参考和视觉轮廓特征等信息。
在 Gazebo 中，上述信息可通过纹理、平面标记或特征板实现[29-30]。与
直接识别起落架外形相比，辅助标识的边界更明确，便于在不同距离和视角下
保持稳定检测，也能使感知模块重点输出身份、中心点和方向差，而不必把复
杂外形识别作为主要任务[30]。
识别码的第一项功能是身份确认。牵引车进入接近区域后，先判断相机视
野中的目标是否与任务编号一致；确认后，再利用标识中心和方向特征计算横
向偏移、距离变化和方向差。这样，图像结果就能转成控制器可以使用的误差
信号，也更具体地体现本文提出的车机互动过程。
为了表述方便，设飞机辅助感知标识的中心点在飞机局部坐标系中的位置
为 𝑃𝑎 ，其在牵引车传感器坐标系中的观测位置为 𝑃𝑠 。则通过图像解算或传感
器映射后，可得到牵引车与目标之间的相对位置关系。若进一步以对接目标中
心作为参考点，则可将识别码的中心点看作前起落架对接区域的等效特征点，
后续所有相对偏差计算均围绕该特征点展开。这样做的目的在于将复杂的机械
接触区域抽象为一个可观测、可追踪、可控制的数学目标，从而显著降低系统
的感知复杂度。


\subsection{感知系统组成与传感器配置}
本课题的感知配置以“视觉识别为主、近距测量为辅”为原则。视觉传感
器负责读取飞机辅助标识并估计目标方向，激光雷达或等效近距传感器负责障
碍物距离、边界补偿和停车阈值判断。该配置没有追求复现真实机场全部感知
复杂度，而是围绕自动接近和最终对接最关键的输入量进行设计。
\subsubsection{视觉识别模块}







视觉模块安装在牵引车前部，镜头朝向前起落架区域。它主要完成识别码
比对、目标中心定位和方向特征提取。
其输出包括目标编号、识别是否有效、图像中心偏移量以及目标方向信
息。
实现方式上，该模块可以使用 RGB 相机或深度相机[29]。若使用深度相
机，还可同步获得目标区域距离，降低单目估距误差。考虑到本文重点是仿真
和控制链路验证，视觉部分只保留“检测—定位—方向估计”的输出接口；后
续进入实物阶段时，可在不改变控制接口的前提下替换为 AprilTag、目标检测
网络或其他算法[29-30]。
\subsubsection{近距测量模块}
近距测量模块用于弥补最终对接阶段中视觉信息的不稳定性。牵引车越靠
近前起落架，目标在图像中越容易被局部放大或遮挡；即使视觉仍能识别标
识，最近距离也未必可靠。因此，本文引入激光雷达、超声波或等效测距模
块，对车体前端与飞机局部结构之间的距离进行校核[18]。
该模块重点输出三类信息：牵引车前端到飞机局部结构的最近距离、非目
标障碍物是否进入安全边界，以及最终接近阶段的停车触发依据。
视觉和近距测量的配合使系统能够随距离变化调整感知重点：中远距离时
先解决“是否看对目标、方向是否大致正确”，近距离时再提高测距、安全边
界和停车条件的优先级。这样更符合牵引车从搜索目标过渡到精确停车的实际
过程。


\subsection{目标身份确认模型}
目标身份确认是自动牵引作业的首要感知任务。只有在确认当前目标为预
定飞机对象之后，牵引车才应进入自动接近与对接控制阶段。否则，即使轨迹
控制精度足够高，也可能因目标对象错误而导致整个作业任务失败。
设任务调度系统下发的目标飞机身份编码为 𝐼𝑡 ，牵引车视觉模块从识别码
中解算得到的观测身份编码为 𝐼𝑜 。则目标身份确认条件可表示为
𝐼𝑜 = 𝐼𝑡        （15）
当上式成立时，判定目标身份匹配；当上式不成立时，系统应保持等待状
态、重新识别或发出异常警告，而不应直接进入对接控制流程。
上述表达式的物理含义非常明确：它并不是一个几何模型，而是一个任务







逻辑判定模型。其作用在于建立从“任务分配”到“目标确认”的第一道闭
环。对于机坪多机位、多牵引任务并行的应用场景而言，这一步尤其重要。若
无身份确认过程，系统只能依赖“看见一个起落架就去接近”的局部策略，这
在真实机场环境中是不可靠的。
𝐶𝐼 𝐶min 𝐶𝐼 ≥ 𝐶min
在仿真实现中，身份确认不宜设计成“一帧识别成功就立即通过”。视觉节点
每次输出识别编号时，同时记录识别置信度；只有目标编号与任务编号一致，
并且置信度连续满足设定阈值时，状态机才把目标有效标志置为真。本文用 表
示当前识别置信度，用 表示确认阈值，则可将附加判据写为 [18,30]。若识别
码短暂丢失、画面遮挡或置信度波动，牵引车仍保持等待或刹停状态，不直接
进入自动接近。这样处理更接近后续程序中的门限判断逻辑，也能避免把偶然
识别结果误作为可靠目标。
\subsection{相对位姿感知模型}
在完成目标身份确认后，牵引车还需进一步获取自身与飞机前起落架区域
之间的相对几何关系。自动接近与对接控制的本质，并不是让牵引车“接近某
个图像目标”，而是让其在物理空间中逐渐消除相对位置误差与方向误差。因
[18,29-30]
此，感知系统必须将图像信息转化为可用于控制的相对位姿量                       。
为此，定义牵引车参考点为 𝑃𝑣 ，飞机辅助感知目标中心为 𝑃𝑎 ，则两者之间
的相对位置误差向量可表示为
𝑒𝑥
𝐞𝑝 = [𝑒 ]                      （16）
𝑦

其中，𝑒𝑥 表示沿牵引车纵向方向的距离误差，𝑒𝑦 表示沿牵引车横向方向
的偏移误差。若从控制意义出发，𝑒𝑥 主要决定前进或减速行为，𝑒𝑦 主要决定
转向修正行为。
同时，定义飞机辅助感知目标区域的参考方向角为 𝜃𝑎 ，牵引车当前航向角
为 𝜃𝑣 ，则相对姿态误差可表示为
𝑒𝜃 = 𝜃𝑎 − 𝜃𝑣                    （17）
式（17）描述的是车辆当前朝向与目标对接方向之间的偏差。车辆即使已
经靠近目标中心，若该角度仍然偏大，后续拖把或连接结构仍可能出现偏接、
卡滞等问题。因此，最终对接不能只看距离，还需要同步检查横向误差和姿态
误差。𝑒𝜃
综合来看，牵引车在自动接近过程中需要实时感知的核心状态量可归纳为






𝐳 = [𝑒𝑥   𝑒𝑦   𝑒𝜃 ]T   （18）
该状态向量构成本章感知环节向控制环节输出的核心接口。换言之，感知
层的任务不是直接控制车辆，而是持续向控制层提供“还差多远、偏了多少、
方向还差多少”的量化描述。
\subsection{近距离态势感知与安全判定}
牵引车进入近距范围后，感知重点会从“找到目标并靠近”转向“确认能
否安全停车”。此时需要关注车体前端与飞机局部结构之间的最近距离、非目
标障碍物是否侵入，以及控制滞后可能造成的提前接触。设前部关键点到飞机
局部结构的最小距离为，安全阈值为，则非接触阶段应始终满足相应距离约束
[5,18]。𝑑min 𝑑𝑠 𝑑min ≥ 𝑑𝑠


最终对接还需要设置停车触发阈值。当剩余纵向距离进入停车范围，并且
横向误差、姿态误差都处于允许范围内时，系统才允许进入低速停止或机械连
接准备状态。设横向误差阈值和姿态误差阈值分别为，则对接有效条件可表示
为𝑑𝑡 𝑒𝑥 ≤ 𝑑𝑡
𝑒𝑦,max 𝑒𝜃,max
∣ 𝑒𝑦 ∣≤ 𝑒𝑦,max
\{                    （19）
∣ 𝑒𝜃 ∣≤ 𝑒𝜃,max
这组条件把距离、横向偏差和方向偏差放在一起判断，避免出现“距离已
经够近但车身还没有对正”的情况。对本文的仿真车而言，只有三个条件同时
满足，状态机才切换到停车保持或连接准备阶段。


\subsection{感知流程设计}
结合上述模型，本文将飞机辅助感知系统的工作流程划分为以下步骤：
\paragraph{} （1）牵引车进入任务区域后，启动视觉识别模块，对前方目标进行搜索；
\paragraph{} （2）检测到飞机辅助感知识别码后，解析目标身份信息，并与任务目标编号进
行匹配；
\paragraph{} （3）若身份匹配成功，则提取目标中心点和方向参考特征，计算相对位置误差
与姿态误差；
\paragraph{} （4）结合近距测量模块输出，对当前局部安全状态进行判断；
\paragraph{} （5）将 𝑒𝑥 、𝑒𝑦 、𝑒𝜃 及安全边界状态发送给控制模块，驱动牵引车完成接近与
对准；







\paragraph{} （6）当对接有效条件满足后，系统进入低速停止或机械连接准备状态。
上述流程对应的是一条从识别到控制的闭环链路。相比只根据图像中心偏
差进行视觉伺服，本文把身份确认和近距安全判断也放入同一流程，更符合机
场牵引这种低速、近距且安全边界要求较严格的任务。
\subsection{控制任务描述与总体思路}
在完成第二章牵引车运动学建模以及前文飞机辅助感知方案设计后，系统
已经具备了两类基础：其一，双桥转向牵引车的平面运动规律可以用数学模型
表示；其二，飞机前起落架区域能够向牵引车提供身份确认结果以及相对位
置、相对姿态等感知信息。在此基础上，本节进一步研究牵引车如何利用这些
信息完成自动接近与最终对接控制。
从实际作业顺序看，牵引车对接前起落架并不是一次性动作。车辆先确认
目标，再驶向预对接区域；进入近距后，控制重点从“快一点靠近”转为“横
向和航向收敛”；最后以很低速度前进，并在距离、姿态和安全条件同时满足
后停车。因此，自动对接控制本质上是任务逻辑、局部目标、误差反馈和停车
判定共同作用的分层闭环问题[25,27-28]。
本文采用“模式管理—局部规划—误差反馈—执行器映射”的控制框架。
感知层持续提供相对误差；模式管理根据距离、障碍物和识别状态选择阶段；
局部规划给出短时域目标点；下层控制器再把速度和转向需求换算成前后桥转
角与四个车轮轮速。这样的结构便于在 ROS 与 Gazebo 中分节点调试，也便于后
续替换感知或控制模块[11-16]。
\subsection{状态变量、控制输入与误差定义}
\subsubsection{车辆状态与控制输入}
依据第二章建立的车辆模型，定义牵引车状态向量为
𝐱 = [𝑥 𝑦 𝜃 𝑣]𝑇              （20）
其中，𝑥 与 𝑦表示牵引车参考点在全局坐标系中的位置坐标，𝜃 表示车辆
航向角，𝑣 表示车辆纵向线速度。
双桥转向牵引车的控制输入定义为
𝐮 = [𝑣        𝛿𝑓   𝛿𝑟 ]𝑇   （21）
或扩展为
𝐮 = [𝑎   𝛿𝑓   𝛿𝑟 ]𝑇   （22）
其中，𝛿𝑓 与 𝛿𝑟 分别为前桥和后桥转角，𝑎 为车辆纵向加速度。前一种形
式适用于直接速度控制实现，后一种形式更适合做速度平滑约束或模型预测控






制扩展。
\subsubsection{感知误差接口}
前文已经将飞机辅助感知结果统一表示为相对误差向量𝐳 = [𝑒𝑥   𝑒𝑦   𝑒𝜃 ]𝑇

——𝑒𝑥 表示牵引车与目标区域之间的纵向距离误差； 𝑒𝑦 表示牵引车相对于目
标轴线的横向偏差； 𝑒𝜃 表示牵引车当前朝向与目标参考方向之间的姿态误差。
在控制意义上，𝑒𝑥 主要决定车辆应继续前进多少距离，𝑒𝑦 主要决定车辆
应向哪一侧修正，𝑒𝜃 则决定车辆当前朝向是否已经满足对接要求。三者共同构
成牵引车自动接近和对准控制的核心反馈量。




\subsection{对接过程分析与分层控制架构}
\subsubsection{对接过程分阶段描述}
为了使控制方法与真实作业过程保持一致，本文将牵引车自动对接过程划
分为以下五个阶段。
\paragraph{} （1）搜索与等待阶段
牵引车进入任务区域后，视觉模块先搜索前起落架附近的识别标识，并与
任务编号比对。若编号不一致或置信度不足，状态机不向下层控制器发布接近
目标，只保持零轮速和转向回中，用于避免误接近。
\paragraph{} （2）自动接近阶段
当目标身份确认有效后，系统开始进入自动接近阶段。该阶段的主要目标
是让牵引车从较远初始位置稳定接近飞机前起落架前方的预对接区域。此时控
制重点不在于最终精确对位，而在于快速减小纵向距离，同时保持可接受的横
向误差和姿态误差。
\paragraph{} （3）姿态对准阶段
当牵引车与目标之间的纵向距离降低到一定阈值后，系统切换到姿态对准
阶段。该阶段中，车辆不再追求较高接近速度，而是优先减小横向误差 𝑒𝑦 和姿
态误差 𝑒𝜃 ，使牵引车中心线逐步与飞机前起落架对接方向重合。
\paragraph{} （4）微动对接阶段
当姿态误差和横向误差已经较小，且车辆进入最终接近区域后，系统进入
微动对接阶段。此时车辆以低速蠕动方式向前推进，控制重点从“大范围修







正”转为“小范围精调”，避免出现冲撞、擦碰或越位现象。
\paragraph{} （5）停车保持阶段
当牵引车与目标区域之间的纵向、横向和姿态误差同时满足终端条件，且
安全距离约束有效时，系统进入停车保持阶段。此时发布零速度命令，并将前
后桥转角回中，保持车辆稳定等待后续机械连接或任务切换。
\subsubsection{分层控制架构}
根据上述阶段分析，本文建立的分层控制架构如图 3-1 所示。
这种架构的优点在于：可以把任务逻辑与连续控制分离，从而降低算法复
杂度，并且更加符合 ROS 节点化实现习惯。




\begin{figure}[htbp]
  \centering
  \includegraphics[width=0.78\textwidth]{pic/fig3-1.png}
  \caption{分层控制架构}
\end{figure}


\subsection{误差动力学模型建立}
为了构造适用于对接过程的控制律，需要首先建立牵引车相对于参考轨迹
或参考目标点的误差动力学模型。
设局部参考状态为𝐪𝑟 = [𝑥𝑟    𝑦𝑟   𝜃𝑟 ]𝑇 。







其中，(𝑥𝑟 , 𝑦𝑟 ) 为局部参考点位置，𝜃𝑟 为参考航向。进一步定义参考速度
和参考角速度分别为 𝑣𝑟 与 𝜔𝑟 。
将全局位置误差变换到车体坐标系，可以得到
𝑒𝑥       cos 𝜃   sin 𝜃 𝑥𝑟 − 𝑥
[𝑒 ] = [               ][      ]
\{ 𝑦       −sin 𝜃 cos 𝜃 𝑦𝑟 − 𝑦              (23)
𝑒𝜃 = 𝜃𝑟 − 𝜃
该式的作用是把全局误差转化为车体系误差，使纵向和横向误差具有明确
物理意义。其中，𝑒𝑥 表示车辆前方还剩多少距离，𝑒𝑦 表示车辆当前偏离对接
轴线的程度。
对上述误差求导，并代入车辆运动学方程，可得误差动力学：
𝑒̇𝑥 = 𝑣𝑟 cos 𝑒𝜃 − 𝑣 + 𝜔𝑒𝑦
\{ 𝑒̇𝑦 = 𝑣𝑟 sin 𝑒𝜃 − 𝜔𝑒𝑥                   (24)
𝑒̇𝜃 = 𝜔𝑟 − 𝜔
上述误差方程表明：
1. 纵向误差 𝑒𝑥 的变化不仅与前进速度 𝑣有关，还受姿态误差和横向运动耦合影
响；
2. 横向误差 𝑒𝑦 的变化主要由角速度控制；
\section{姿态误差 𝑒𝜃 的变化取决于参考角速度与实际角速度的差异。}
这说明自动对接控制本质上是一个非线性耦合误差调节问题，不能简单把纵
向和横向完全独立看待[19,27-28]。
\subsection{局部参考轨迹生成方法}
\subsubsection{预对接点与对接点构造}
为了避免车辆直接朝最终机械连接位置冲击，本文采用两级目标点设计，即
“预对接点”和“对接点”。
设飞机辅助感知系统给出的目标中心点为
𝑥𝑎
𝐩𝑎 = [𝑦 ]                (25)
𝑎

目标方向为 𝜃𝑎 。则预对接点定义为
cos 𝜃𝑎
𝐩𝑝 = 𝐩𝑎 − 𝑑pre [          ]    （26）
sin 𝜃𝑎
𝑑pre 为预留接近距离。该点位于最终目标点前方一定距离处，主要用于完
成姿态调整。最终对接点定义为
cos 𝜃𝑎
𝐩𝑑 = 𝐩𝑎 − 𝑑hook [          ]   （27）
sin 𝜃𝑎








𝑑hook 该点为机械连接前的停车位置，主要用于微动阶段。若车辆从较远位
置直接追踪最终点，接近后段容易出现横向修正不足或转向过猛；设置预对接
点后，车辆可以先把姿态调到接近目标方向，再以较低速度进入最后一段。
\subsection{参考速度调度策略}
在机场牵引对接作业中，恒定速度控制并不合适。当路径弯曲程度较大
时，高速运行会导致横向误差难以收敛； 而当障碍物距离较小时，必须主动减
速保证停车裕度； 当进入微动阶段时，也需要明显降低速度以提高最终对位精
度。 因此，参照相关文献资料，认为可以采用如下速度调度公式：
𝑎𝑦,max
𝑣𝑟 = min (𝑣stage , √∣𝜅 ∣+𝜀 , 𝑘𝑑 [𝑑min − 𝑑𝑠 ]+ )       (28）
𝑟   𝜅


其中：𝑣stage 为当前控制模式允许的最大速度； 𝑎𝑦,max 为允许的最大横向加
速度； 𝜅𝑟 为参考轨迹曲率； 𝑑min 为前方最近障碍距离； 𝑑𝑠 为安全距离阈值；
𝑘𝑑 为距离—速度缩放系数[25,27-28]。
\subsection{误差反馈控制律设计}
\subsubsection{纵向速度控制律}
纵向控制律设计为
𝑣𝑑 = sat [0,𝑣max ] (𝑣𝑟 cos 𝑒𝜃 + 𝑘𝑥 𝑒𝑥 )           (29）

其中，纵向误差增益用于调节距离误差对速度的影响，饱和函数用于限制
速度输出。𝑘𝑥 sat
该控制律的作用可以理解为两步：姿态偏差较大时先压低前进速度，避免
未对正就继续推进；纵向误差较大时再给出适当推进速度，使车辆逐步靠近目
标。𝑣𝑟 cos 𝑒𝜃 𝑘𝑥 𝑒𝑥
因此，车辆距离目标较远时可以保持一定接近速度；一旦航向偏差变大，
速度会被主动压低，为后续转向修正留下空间[19]。
\subsubsection{角速度控制律}
角速度控制律设计为

𝜔𝑑 = sat [−𝜔max ,𝜔max ] (𝜔𝑟 + 𝑣𝑟 (𝑘𝑦 𝑒𝑦 + 𝑘𝜃 sin 𝑒𝜃 ))   （30）

其中，𝑘𝑦 和 𝑘𝜃 分别为横向误差增益和姿态误差增益。该控制律的物理意义如
下[17,26-28]：
𝜔𝑟 ：参考轨迹自身的前馈角速度，用于保证车辆沿局部参考曲线运行；
𝑣𝑟 𝑘𝑦 𝑒𝑦 ：根据横向偏差调节转向，使车辆向目标轴线靠拢；







𝑣𝑟 𝑘𝜃 sin 𝑒𝜃 ：根据姿态误差修正车辆方向。
之所以采用 sin 𝑒𝜃 而不是直接用 𝑒𝜃 ，是因为前者在大角度误差情况下更平
滑，且在小角度时近似等于 𝑒𝜃 ，便于后续稳定性分析。
\subsubsection{对接过程中的控制行为分析}
结合车辆对接全过程，上述控制律在不同阶段表现如下：
\paragraph{} （1）自动接近阶段：
此时 𝑒𝑥 较大，速度项 𝑘𝑥 𝑒𝑥 占主导，车辆以较快速度向预对接点推进。
\paragraph{} （2）姿态对准阶段
随着 𝑒𝑥 逐步减小，横向误差 𝑒𝑦 和姿态误差 𝑒𝜃 开始主导控制器输出，角速度命
令增大，车辆重点完成方向修正。
\paragraph{} （3）微动对接阶段：速度调度将线速度限制在较低范围，角速度输出也随之变
小，控制器只做小幅修正，减少末端摆动和越位。𝑣𝑟
由此可见，上述控制律并不是孤立给出的公式，而是与“接近—对准—微
动”的动作顺序相对应，便于在 Gazebo 中观察不同阶段的车辆响应。
\subsection{四轮转向映射与约束处理}
\subsubsection{曲率计算}
由期望线速度和角速度得到期望曲率
𝜔
𝜅𝑑 = max (𝑣𝑑 ,𝜀 )                  （31）
𝑑 𝑣


其中, 𝜀𝑣 为防止低速时分母趋近于零的安全常数。


\subsubsection{双桥反向转向模式}
在自动接近阶段，如果需要较大机动能力，则采用双桥反向转向模式[19,28]：
𝐿𝜅𝑑
𝛿𝑓⋆ = sat (arctan         , −𝛿
𝑓,max     , 𝛿𝑓,max )
\{                     2                            (32)
⋆     ⋆
𝛿𝑟 = −𝛿𝑓

该模式可以显著减小等效转弯半径，适合从较远位置快速接近飞机前起落
架区域。
\subsubsection{前桥单转向模式}
在姿态对准和微动阶段，为了提高控制平稳性，采用前桥单转向模式：

𝛿 ⋆ = sat (arctan (𝐿𝜅𝑑 ), −𝛿𝑓,max , 𝛿𝑓,max )
\{ 𝑓                                                (33)
𝛿𝑟⋆ = 0








该模式虽然机动性不如双桥反向转向，但执行器协调更简单，转向响应更
平顺，适合最终精调。
\subsubsection{速度与转角限幅}
为了保证控制命令可执行，还需对速度与转角进行约束处理：
0 ≤ 𝑣𝑑 ≤ 𝑣max
⋆
\{ ∣ 𝛿𝑓 ∣≤ 𝛿𝑓,max                                               (34)
∣ 𝛿𝑟⋆ ∣≤ 𝛿𝑟,max
同时对速度变化率和转角变化率限制为
∣ 𝑣𝑑 (𝑘) − 𝑣𝑑 (𝑘 − 1) ∣≤ 𝑎max 𝑇𝑠
\{                            ̇ 𝑇𝑠 , 𝑖 ∈ \{𝑓, 𝑟\}                      (35)
∣ 𝛿𝑖 (𝑘) − 𝛿𝑖 (𝑘 − 1) ∣≤ 𝛿max
这样可以避免由于离散控制导致的命令突变，从而减小 Gazebo 中常见的抖
振与过冲现象。
\subsection{终端停车判定与模式切换}
\subsubsection{终端停车集合}
在最终微动阶段，当牵引车同时满足以下条件时，可判定已具备停车条
件：
∣ 𝑒𝑥 ∣≤ 𝑑𝑡
∣ 𝑒𝑦 ∣≤ 𝑒𝑦,max
\paragraph{} (36)
∣ 𝑒𝜃 ∣≤ 𝑒𝜃,max
\{ 𝑑min ≥ 𝑑𝑠
将其定义为终端停车集合

Ω𝑡 = \{∣ 𝑒𝑥 ∣≤ 𝑑𝑡 , ∣ 𝑒𝑦 ∣≤ 𝑒𝑦,max , ∣ 𝑒𝜃 ∣≤ 𝑒𝜃,max , 𝑑min ≥ 𝑑𝑠 \}   (37)

与单纯依靠纵向距离判定停车不同，这种集合式判定可以避免“距离够近
但方向没对准”的错误停车问题。
\subsubsection{模式切换规则}
系统模式切换规则如下：
搜索-自动接近：识别码身份确认成功且置信度满足阈值；
自动接近-姿态对准：𝑒𝑥 进入近距阈值；
姿态对准-微动对接：𝑒𝑥 足够小且 𝑒𝜃 已较小；
任意模式-中止：目标丢失或最近距离低于中止阈值；
微动对接-停车保持：进入终端集合 Ω𝑡 。
这意味着控制系统并不是依赖一条连续控制律从头走到尾，而是通过“阶
段切换 + 连续控制”的混合方式实现更符合工程需求的动作组织。





